\documentclass[11pt]{article}
\usepackage[a4paper,margin=1in]{geometry}
\usepackage[T1]{fontenc}
\usepackage[utf8]{inputenc}
\usepackage{times}
\usepackage{parskip}
\usepackage[hidelinks]{hyperref}

\begin{document}

\begin{flushright}
October 21, 2026
\end{flushright}

\vspace{0.5em}

The Editor-in-Chief \\
\textit{Computer Modeling in Engineering \& Sciences} (CMES) \\
Tech Science Press

\vspace{1em}

Dear Editor,

\vspace{0.5em}

\textbf{Re: Submission of a new manuscript for consideration in CMES}

We are pleased to submit our manuscript entitled \textit{``Fixed Points and Certified
Boundedness of Unrolled Born/Contrast-Source Inverse Scattering Networks''} for
consideration as a Research Article in \textit{Computer Modeling in Engineering \&
Sciences}. The corresponding author is Longgang Wang (Luoyang Institute of Science
and Technology, Luoyang 471000, Henan, China; email: longgangwang@163.com; ORCID:
0009-0008-1394-7953).

\subsection*{Originality and author approval}

This manuscript is original and has not been published previously. It is not under
consideration for publication elsewhere, in whole or in part. All authors have read
the final manuscript and approved its submission to the journal. They accept full
responsibility for the manuscript's delivery and contents. The work reports no human
or animal subjects. There are no copyright or disclosure issues associated with the
manuscript. All authors declare that they have no conflicts of interest regarding
the submitted manuscript.

\subsection*{What the paper contributes}

Deep unrolling of classical iterative solvers is a leading design principle for
electromagnetic inverse scattering, and unrolled architectures reported in the
literature outperform their classical ancestors empirically. What has been missing is a
layer-level account of \emph{why} such a network is stable, when it is provably not,
and what can be certified about it. This paper supplies that account. Its four
contributions, which mirror the manuscript's contribution list, are:

\begin{enumerate}
\item \textbf{Layer-level operator form.} The discrete Lippmann--Schwinger system is
written as a fixed-point iteration of a physics operator $T^{\mathrm{phys}}$ on the
state $z=(\chi,w)$; an unrolled network is exactly a cascade of layers
$T_k=(1-\theta_k)I+\theta_k T^{\mathrm{phys}}_{\alpha_k,\gamma_k}+H_k$ with $H_k$ a
learned correction, and the fixed points of the trained map are characterized exactly.

\item \textbf{Impossibility of global contraction, in threshold form.} For the faithful
Born/contrast-source composite, the bilinear coupling $\chi\odot E$ makes one block of
the layer Jacobian \emph{exactly} $\theta\alpha\,\mathrm{diag}(E)$. Because that block
belongs to the second block row, which the learned correction does not touch, the bound
$\sigma_{\max}(DT_k)\ge\theta\alpha\mathcal{E}>1$ holds \emph{for every correction and
every Lipschitz constant} above an explicit step-size threshold. Global Euclidean
contraction is therefore unattainable on the measured step-size grid, and stability
guarantees must be layered rather than global.

\item \textbf{Layer-wise certification instead of per-layer contraction.} We prove
what \emph{is} available: an a-priori field-subsystem contraction
($\beta_0=\chi_{\max}\|\mathbf{G}_b\|_2<1$) that confines the rollout to an explicit
compact set at arbitrary depth under a nonexpansive physics-prior projection, so
infinite-depth extrapolation cannot diverge; an exact finite-depth noise accounting,
including the non-normal factor $\kappa_2(V)$ that a naive $\Lambda/(1-\rho)$ omits;
and an explicit constructive bound on the data-Lipschitz constant. The geometric bounds
hold only \emph{conditionally}, and we demonstrate on our own measurements where their
hypotheses fail. The certificates are recomputed at the operator level on the published
Institut Fresnel geometry, where the weak-scattering condition fails while the
non-contractivity bound tightens, and a reconstruction from measured fields is reported
for nine phantom-frequency pairs without the trained weights.

\item \textbf{The price of certification, measured.} Pre-registered ablations over 800
unrolled layers isolate which ingredient delivers which guarantee: removing the
projection from frozen weights makes the unconstrained network explode, and retraining
without it restores empirical boundedness but pushes the measured contrast bound past
the weak-scattering threshold, so no a-priori certificate survives. The amplitude
budget buys the a-priori constants at a measurable accuracy cost, and bounding the
correction trunk's input is what keeps the state inside the admissible box.
\end{enumerate}

The paper also reports a per-layer cost audit and an identical-seed reproducibility
floor, so that each numerical claim can be traced to a script and an archived artifact.

\subsection*{Why this journal}

CMES publishes research of the highest quality in computational modeling and engineering
sciences. Our contribution is squarely within its scope: methodologically, it combines
operator theory of the discrete Lippmann--Schwinger formulation with contraction-mapping
and spectral analysis to obtain fixed-point and boundedness certificates for unrolled
inverse-scattering networks, and the certifying constants are computed on the 2-D TM
method-of-moments scattering operator, an inverse/imaging problem in the journal's stated
application domain. The framework is a reusable template that separates what can be
certified from what is only measured.

\subsection*{Data and code availability}

All numerical instruments, experiment drivers, audit scripts, and a frozen snapshot with
a standard \texttt{md5sum -c} manifest are listed in the manuscript's Data and Code
Availability section. A public archive URL or persistent identifier will accompany the
final version.

\subsection*{Suggested reviewers}

We are glad to suggest potential reviewers for this manuscript upon request.

We thank you for considering our manuscript and look forward to your response.

\vspace{1.5em}

Sincerely,

\vspace{1.5em}

Longgang Wang \\
on behalf of all authors \\
Luoyang Institute of Science and Technology \\
Luoyang 471000, Henan, China \\
Email: longgangwang@163.com

\end{document}
